% =============================================================================
\chapter{Painel interativo}\label{cap:dashboard}
% =============================================================================

O painel interativo, implementado em \code{app/dashboard.py} com Streamlit
\cite{streamlit2026} e Plotly \cite{plotly2015}, é o produto de visualização do
\textit{Checkpoint}~3. Ele permite que a mesa de execução explore os indicadores
das perguntas P1 a P3 com filtros por dia, janela horária, fuso e faixa de
tamanho, e consulte os resultados da modelagem.

\section{Arquitetura e desempenho}\label{sec:dasharquitetura}

O painel não acessa a base bruta: ele consome somente os Parquets de
\code{data/processed/}, por meio das funções de \code{src/data_loader.py} e
\code{src/modeling.py}. Se os arquivos não existirem, exibe uma mensagem pedindo
que o ETL seja executado antes. Para manter a interação fluida, as operações
custosas são memorizadas com o decorador \code{st.cache_data} e executadas uma
única vez por sessão do servidor:

\begin{alineas}
  \item leitura das barras de 1 minuto;
  \item resumos centrais: resumo diário, distribuição de negócios por dia, hora e
        faixa de tamanho (com \textit{cache} adicional em
        \code{trade_size_hourly.parquet}) e agregação das ordens agressoras por
        dia, hora e faixa, incluindo a contagem e o volume das ordens de
        \usd~100 mil ou mais;
  \item validação \textit{walk-forward} completa em barras de 1 minuto, para as
        tarefas de volatilidade e direção.
\end{alineas}

Os filtros são aplicados depois, sobre esses resumos já agregados, o que evita
reprocessar os 27,9 milhões de negócios a cada clique.

\section{Filtros}\label{sec:filtros}

\begin{quadro}[htbp]
\caption{\label{qua:filtros}Filtros da barra lateral do painel}
\begin{corpotabela}
\begin{tabularx}{\textwidth}{|L{3.4cm}|L{4.2cm}|Y|}
\hline
\textbf{Filtro} & \textbf{Opções} & \textbf{Efeito} \\ \hline
Fuso horário & UTC; BRT (UTC$-3$) & Desloca o eixo horário de todos os perfis, barras e mapas de calor \\ \hline
Recorte rápido de dias & Semana completa; apenas dias de estresse (21, 24 e 25/02); apenas dias calmos (19, 20, 22 e 23/02); seleção manual & Pré-seleciona os dias do filtro seguinte \\ \hline
Dias incluídos & Cada dia, rotulado com o tipo (por exemplo, ``ter 25/02 $\cdot$ estresse'') & Restringe todos os gráficos e indicadores aos dias escolhidos \\ \hline
Janela horária & Intervalo de 0h a 23h no fuso escolhido & Restringe barras, perfis e agregações às horas escolhidas \\ \hline
Faixas de tamanho & As cinco faixas de \code{size_bucket} & Restringe as análises de concentração por tamanho (P2) \\ \hline
\end{tabularx}
\end{corpotabela}
\fonteautores[Elaborado]
\end{quadro}

Duas decisões de projeto devem ser conhecidas por quem usa o painel. Primeiro, o
filtro de dias usa a data em UTC; ao escolher BRT, muda o rótulo das horas, mas
cada dia continua delimitado pela meia-noite UTC. Segundo, as tabelas de
modelagem da aba P3 resumem a validação \textit{walk-forward} da semana inteira e
não respondem aos filtros, porque o treino depende da sequência completa de dias.

\section{Indicadores e abas}\label{sec:abas}

O topo do painel exibe o título, um resumo do filtro ativo e cinco indicadores
recalculados a cada mudança de filtro. Abaixo deles, quatro abas organizam o
conteúdo segundo a narrativa do projeto. A \autoref{fig:wireframe} mostra um
esquema da interface, e o \autoref{qua:abas} detalha o conteúdo de cada aba.

\begin{figure}[htbp]
\caption{\label{fig:wireframe}Esquema da interface do painel interativo}
\begin{center}
\begin{tikzpicture}[font=\scriptsize, x=1cm, y=1cm]
  % moldura
  \draw[cinza, line width=0.6pt, rounded corners=2pt] (0,0) rectangle (15.8,9.6);
  % barra lateral
  \fill[cinzaclaro] (0.05,0.05) rectangle (3.9,9.55);
  \node[anchor=north west, font=\footnotesize\bfseries] at (0.25,9.4) {Filtros da Mesa};
  \foreach \y/\t in {8.6/Fuso horário, 7.7/Recorte rápido de dias, 6.8/Dias incluídos, 5.9/Janela horária, 5.0/Faixas de tamanho}{
    \node[anchor=west, text=tinta] at (0.25,\y+0.15) {\t};
    \draw[cinza, fill=white, rounded corners=1.5pt] (0.25,\y-0.45) rectangle (3.7,\y-0.05);
  }
  \node[anchor=north west, text width=3.4cm, align=left, text=tinta, font=\tiny] at (0.25,3.9)
    {Fonte: Binance \textit{spot} BTC/USDT via Hugging Face (\texttt{c90f6087}), reconciliada 100\,\% com os candles oficiais.};
  % título e legenda
  \node[anchor=north west, font=\small\bfseries] at (4.2,9.4) {BTC/USDT -- Microestrutura da Calmaria ao Estresse};
  \node[anchor=north west, text=tinta] at (4.2,8.85) {Filtro atual: dias, janela horária, fuso e número de minutos};
  % KPIs
  \foreach \i/\t in {0/Volume, 1/Negócios, 2/Amplitude, 3/Faixa de preço, 4/Saldo comprador}{
    \pgfmathsetmacro{\x}{4.2+\i*2.3}
    \draw[azul, fill=azulclaro, rounded corners=2pt] (\x,7.25) rectangle (\x+2.15,8.3);
    \node[anchor=north, font=\tiny, text=tinta] at (\x+1.075,8.22) {\t};
    \node[font=\footnotesize\bfseries] at (\x+1.075,7.6) {valor};
  }
  % abas
  \foreach \i/\t in {0/Calmaria ao estresse, 1/P1 $\cdot$ Liquidez, 2/P2 $\cdot$ Estresse, 3/P3 $\cdot$ OFI e modelos}{
    \pgfmathsetmacro{\x}{4.2+\i*2.88}
    \ifnum\i=0
      \draw[laranja, fill=laranjaclaro] (\x,6.45) rectangle (\x+2.8,6.95);
    \else
      \draw[cinza, fill=white] (\x,6.45) rectangle (\x+2.8,6.95);
    \fi
    \node[font=\tiny] at (\x+1.4,6.7) {\t};
  }
  % conteúdo
  \draw[cinza, fill=white] (4.2,3.0) rectangle (15.6,6.3);
  \node[text=tinta, align=center] at (9.9,4.65) {Preço de fechamento por hora e volume horário colorido pelo tipo de dia\\(gráficos interativos Plotly)};
  \foreach \i/\t in {0/{1. Janela TWAP/VWAP (P1)}, 1/{2. Limites e capacidade (P2)}, 2/{3. Uso realista de sinais (P3)}}{
    \pgfmathsetmacro{\x}{4.2+\i*3.85}
    \draw[azul, fill=azulclaro!60, rounded corners=2pt] (\x,0.3) rectangle (\x+3.7,2.75);
    \node[anchor=north, text width=3.4cm, align=center, font=\scriptsize\bfseries] at (\x+1.85,2.65) {\t};
    \node[anchor=north, text width=3.4cm, align=center, text=tinta] at (\x+1.85,2.05) {recomendação para a mesa baseada nos achados};
  }
\end{tikzpicture}
\end{center}
\fonte{Elaborada pelos autores (2026). Esquema ilustrativo da disposição dos elementos de \texttt{app/dashboard.py}, não uma captura de tela.}
\end{figure}

\begin{quadro}[htbp]
\caption{\label{qua:abas}Indicadores e conteúdo das abas do painel}
\begin{corpotabela}
\begin{tabularx}{\textwidth}{|L{3.3cm}|Y|}
\hline
\textbf{Elemento} & \textbf{Conteúdo} \\ \hline
Indicadores do topo & Volume negociado (\usd~bilhões); negócios executados (milhões); amplitude mediana de 1 minuto (bps); faixa de preço (mínima a máxima); saldo comprador líquido (\% do volume) \\ \hline
Aba ``Da Calmaria ao Estresse'' & Preço de fechamento por hora e volume horário colorido pelo tipo de dia; plano de ação da mesa em três cartões: janela TWAP/VWAP (P1), limites e capacidade em estresse (P2) e uso realista de sinais preditivos (P3) \\ \hline
Aba ``P1 $\cdot$ Liquidez Intradiária'' & Medianas por hora de volume, negócios e amplitude, com as três horas de maior volume em destaque; mapa de calor de volume por dia e hora; dispersão entre volume e amplitude por hora, com o $\rho$ de Spearman no título \\ \hline
Aba ``P2 $\cdot$ Estresse \& Faixas de Tamanho'' & Volume e volatilidade realizada por dia, coloridos pelo tipo de dia; tabela de médias por regime; participação de cada faixa na contagem e no valor dos negócios; número de ordens agressoras de \usd~100 mil ou mais por dia \\ \hline
Aba ``P3 $\cdot$ Fluxo (OFI) \& Modelagem'' & Saldo líquido diário do fluxo agressor; correlação de Pearson entre $\mathrm{OFI}_t$ e $r_{t+k}$ para $k$ de 0 a 10, com a faixa de ruído de 95\,\%; tabelas-resumo da validação \textit{walk-forward} para volatilidade e direção \\ \hline
\end{tabularx}
\end{corpotabela}
\fonteautores[Elaborado]
\end{quadro}

\section{Execução}\label{sec:dashexecucao}

Com os Parquets gerados (\autoref{sec:comoexecutar}), o painel é iniciado a
partir da raiz do repositório com o comando abaixo e abre no navegador, por
padrão em \url{http://localhost:8501}.

\begin{lstlisting}[language=bash]
streamlit run app/dashboard.py
\end{lstlisting}

A primeira carga é a mais lenta, porque executa a validação
\textit{walk-forward} e grava o \textit{cache} de faixas de tamanho; as interações
seguintes reutilizam os resultados memorizados. O teste automatizado
\code{test_streamlit_dashboard_renders_without_errors} executa o painel sem
navegador e verifica que ele renderiza sem exceções, com cinco indicadores e
quatro abas.
